\documentclass[12pt]{article}
%% BUGS:
%% - The \begin{problem} ... \end{problem} is only used in exams??

\usepackage{url, graphicx, epstopdf}

% page layout
\usepackage[letterpaper, textheight=9.5in, textwidth=5.5in]{geometry}
\usepackage{setspace}
\setstretch{1.08}
\pagestyle{plain}
\sloppy\sloppypar\raggedbottom\frenchspacing\thispagestyle{empty}

% problem set header
\newcommand{\psname}{Problem Set}
\newcommand{\startps}[1]{\section*{Computational Physics --- \psname~{#1}}}

% problem formatting
\usepackage{amsthm}
\theoremstyle{definition}
\newcommand{\problemname}{Problem}
\newtheorem{problem}{\problemname}
\newcommand{\probref}[1]{\problemname~\ref{#1}}
\theoremstyle{remark}
\newcommand{\notename}{Note}
\newtheorem{note}{\notename}[problem]
\newcommand{\noteref}[1]{\textit{\notename}~\ref{#1}}
\newcounter{subpart}[problem]
\newcommand{\subpart}{\par\addvspace{1ex}\refstepcounter{subpart}\textsl{(\alph{subpart})}~}

% words
\newcommand{\foreign}[1]{\textsl{#1}}
\newcommand{\vs}{\foreign{vs}}

% code
\newcommand{\code}[1]{\texttt{\detokenize{#1}}}
\usepackage{xcolor}
% Define custom colors for code highlighting
\definecolor{codegreen}{rgb}{0,0.6,0}
\definecolor{codegray}{rgb}{0.5,0.5,0.5}
\definecolor{codepurple}{rgb}{0.58,0,0.82}
\definecolor{backcolour}{rgb}{0.99,0.99,0.99}
\usepackage{listings}
\lstdefinestyle{pythonstyle}{
    backgroundcolor=\color{backcolour},   
    commentstyle=\color{codegreen},
    keywordstyle=\color{magenta},
    numberstyle=\tiny\color{codegray},
    stringstyle=\color{codepurple},
    tabsize=4
}
\lstset{
    basicstyle=\small\ttfamily,
    columns=fixed,
    breaklines=false,
    keepspaces=true,
    showspaces=false,
    showstringspaces=false,
    showtabs=false,
    showlines=*,
    style=pythonstyle,
    language=Python,
    literate={~}{{\url{~}}}1
}

% math
\usepackage{amsmath}
\newcommand{\dd}{\mathrm{d}}
\newcommand{\e}{\mathrm{e}}

% primary units
\newcommand{\rad}{\mathrm{rad}}
\newcommand{\kg}{\mathrm{kg}}
\newcommand{\m}{\mathrm{m}}
\newcommand{\s}{\mathrm{s}}

% secondary units
\renewcommand{\deg}{\mathrm{deg}}
\newcommand{\km}{\mathrm{km}}
\newcommand{\cm}{\mathrm{cm}}
\newcommand{\mm}{\mathrm{mm}}
\newcommand{\ft}{\mathrm{ft}}
\newcommand{\mi}{\mathrm{mi}}
\newcommand{\AU}{\mathrm{AU}}
\newcommand{\ns}{\mathrm{ns}}
\newcommand{\h}{\mathrm{h}}
\newcommand{\yr}{\mathrm{yr}}
\newcommand{\N}{\mathrm{N}}
\newcommand{\J}{\mathrm{J}}
\newcommand{\eV}{\mathrm{eV}}
\newcommand{\W}{\mathrm{W}}
\newcommand{\Pa}{\mathrm{Pa}}

% derived units
\newcommand{\mps}{\m\,\s^{-1}}
\newcommand{\mph}{\mi\,\h^{-1}}
\newcommand{\mpss}{\m\,\s^{-2}}

% random units
\newcommand{\au}{\mathrm{a.u.}}


\begin{document}

\startps{1}
Due Friday 2026 September 11 at 17:00 on Brightspace.
Figure out a way to make your solution be one big PDF file with figures and text; maybe use \LaTeX?
Be sure to explicitly give credit where it is due, and as specifically as possible, for all your problems.

\begin{problem}
  \subpart What is the precision of a floating-point number on the computer?
  Experiment by adding very small numbers to the number 1.0 until you see no change.
  That is, what is the smallest $\delta$ such that \emph{on the computer} $[1+\delta]>1$?
  Is your answer for \code{float} or \code{double}?

  \subpart Now what is the largest $\Delta$ such that \emph{on the computer} $[1+\Delta]>\Delta$?

  \subpart What is the largest $\Delta_1$ such that $\exp(-\Delta_1)>0$?
  What about the largest $\Delta_2$ such that $\exp(\Delta_2)$ gives a floating-point answer?
  And what are the corresponding values of $\exp(-\Delta_1)$ and $\exp(\Delta_2)$?
\end{problem}

\begin{problem}
  \subpart How far does a muon travel during its mean lifetime?
  This depends on how fast it is moving.
  Look up the lifetime of the muon and make a plot of the distance traveled by the muon during that lifetime as a function of the speed it is going.
  Make a plot of distance (on the vertical axis) \vs{} dimensionless speed $\beta$ and make a plot \vs{} Lorentz factor $\gamma$.
  Plot a labeled horizontal line showing the naive maximum distance (that is, the mean lifetime times the speed of light $c$) and a labeled horizontal line showing the scale height of the atmosphere ($10\,\km$).
  Why is that latter distance relevant?

  \subpart Now change to plotting distance \vs{} $[1-\beta]$ and vs $\gamma$, but make your plots log--log.
  Add a labeled horizontal line showing the distance from the Earth to the Sun.
  Make sure your range of $[1-\beta]$ and $\gamma$ reaches to (or, better, slightly beyond) the Solar distance!
\end{problem}

\begin{note}
  You have finished this problem when you have produced four plots, all with good labels and clear content.
  Bonus points if these four plots are produced as four panels in a $2\times 2$ plot matrix.
\end{note}

\begin{note}
  The speed $\beta$ is a bad way to parameterize high speeds; $\gamma$ is better.
  How do you parameterize both low speeds and high speeds in a way that works okay, computationally?
\end{note}

\begin{note}
  The point of this problem is to make figures that can be understood by someone, or used in a textbook, say.
  Do your figures meet these criteria?
\end{note}

\begin{problem}
  Here's a function of two angles, $\theta,\varpi$:
  \begin{align}
    f(\theta,\alpha) &= \sin(2\,\theta + 1)\,\sin(3\,\varpi - \frac{\pi}{4}) ~.
  \end{align}
  \subpart Plot this function with a sensible colormap on the square $0<\theta<2\pi$, $0<\varpi<2\pi$.
  Make sure the plot range is exactly this square range and that the two axes are properly labeled.
  When you do this, you will have to decide on a size (number of pixels) for the underlying image.
  Make your underlying image exactly 18 pixels in the $\theta$ direction and 30 pixels in the $\alpha$ direction.
  Overplot thin or faint vertical and horizontal lines showing where you expect to see maxima of this function.

  \subpart Now plot this function $f(\theta,\alpha)$ taken to the fifth power.
  The peaks should get sharper, and should be precisely located in the correct places.
  Or should they?
  When you make all your plots, make sure you use \code{interpolation="nearest"} or equivalent.

  \subpart Now make the two plots again, but use a grid that is 521 pixels in the $\theta$ direction and 601 pixels in the $\varpi$ direction.
  Do you feel better about your plots now? Or not?
\end{problem}

\begin{note}
  When you make a finite pixelized image, you have to think about where the left edge of each pixel is, where the right edge of each pixel is, and where the center of each pixel is.
  Think hard about this and choose wisely!
  I am very specific about the plotting range; I want it to be exactly $0<\theta<2\pi$ and $0<\varpi<2\pi$.
\end{note}

\begin{problem}
  Write down a few tests of the \code{sin()} and \code{cos()} functions, based on old-school trig identities, or even simpler facts about angles.
  Choose tests that depend only on one angle $\theta$.
  Each of these tests can be written to produce a residual; for example, the fact that $\sin^2\theta+\cos^2\theta=1$ can be written in residual form as
  \begin{align}
    \delta &= \sin^2\theta + \cos^2\theta - 1 ~.
  \end{align}
  Plot, for each of your tests, the value of $\delta$ as a function of $\log_{10}(\theta)$.
  Give your plots good labels and title or label each plot with the test that it is executing.
  Extend your plots far enough in $\theta$ that you start to see issues.
  Or do you ever see issues?
\end{problem}

\begin{note}
  You might look at problems that involve adding angles or dividing angles, to see the numerical issues most clearly.
  For example, a dumb test is $\sin(\theta + 2\pi) = \sin(\theta)$.
\end{note}

\begin{note}
  Are your plots informative about, or useful for, anything?
  Is it more useful to plot the residual $\delta$ or (when it is non-zero) the logarithm of it?
\end{note}

\end{document}
